\documentclass[11pt]{article}

\usepackage[margin=1in]{geometry}
\usepackage{amsmath,amssymb,amsthm}
\usepackage{graphicx}
\usepackage[hidelinks]{hyperref}

\newtheorem{theorem}{Theorem}
\newtheorem{lemma}{Lemma}
\newtheorem{remark}{Remark}

\newcommand{\Lip}{\operatorname{Lip}}
\newcommand{\Free}{\mathcal F}
\newcommand{\Ball}{B}
\newcommand{\norm}[1]{\left\lVert #1 \right\rVert}

\title{Finite Metric Spaces in a Lipschitz-Free Uniform Property \texorpdfstring{$(S)$}{(S)} Question}
\author{Automated Math Discovery Smoke Test}
\date{June 7, 2026}

\begin{document}
\maketitle

\section*{Source and Open Question}

This packet concerns the Lipschitz-free item in the open-problem section
of:

\begin{quote}
William B. Johnson and Tomasz Kania, \emph{Uniform Property (S)},
arXiv:2602.09106.
\end{quote}

The source location is page 19 of the paper.

\begin{center}
\includegraphics[width=\textwidth]{figures/open_problem_crop.png}
\end{center}

The paper defines property $(S)$ as follows: a normed space $X$ has
property $(S)$ if for every distinct $x,y\in S_X$ and every $a>0$ there is
$z\in X$ with $\norm{z}<a$ such that
\[
  \norm{x+z}\neq \norm{y+z}.
\]
Uniform property $(S)$ is the corresponding quantitative strengthening
using the modulus
\[
  U_X(d;a)=
  \inf_{\substack{x,y\in S_X\\ \norm{x-y}\ge d}}
  \sup_{\norm{z}\le a}
  \left|\norm{x+z}-\norm{y+z}\right|.
\]
Uniform property $(S)$ implies property $(S)$.

\section*{Candidate Partial Result}

\begin{theorem}
Let $M$ be a finite pointed metric space. If $|M|\ge 3$, then the
Lipschitz-free Banach space $\Free(M)$ fails property $(S)$, and therefore
fails uniform property $(S)$.

If $|M|=2$, then $\Free(M)$ has uniform property $(S)$. If $|M|=1$, the
space $\Free(M)$ is the zero space, so the usual unit-sphere formulation
is vacuous.
\end{theorem}

\begin{proof}
We use the standard description of finite Lipschitz-free spaces:
\[
  \Ball_{\Free(M)}
  =
  \operatorname{conv}
  \left\{
  \frac{\delta_u-\delta_v}{d(u,v)}:
  u,v\in M,\ u\neq v
  \right\}.
\]
Thus, when $M$ is finite, $\Ball_{\Free(M)}$ is a centrally symmetric
polytope. Also $\dim \Free(M)=|M|-1$.

Assume first that $|M|\ge 3$. Then $\dim\Free(M)\ge 2$, so the polytope
$\Ball_{\Free(M)}$ has a facet $F$ of positive dimension. Choose two
distinct points $x,y$ in the relative interior of $F$. Let $\varphi$ be a
norm-one functional exposing this facet, so that
\[
  F=\{u\in \Ball_{\Free(M)}:\varphi(u)=1\}.
\]

We claim that the norm is locally equal to $\varphi$ near both $x$ and
$y$. Indeed, because the ball is a polytope, there are finitely many
norm-one functionals $\varphi_1,\ldots,\varphi_N$ such that
\[
  \norm{u}=\max_{1\le j\le N}\varphi_j(u)
  \quad (u\in \Free(M)),
\]
and one of them is $\varphi$. Since $x$ and $y$ lie in the relative
interior of the facet exposed by $\varphi$, all the other facet
functionals are strictly smaller than $\varphi$ at both $x$ and $y$.
By finiteness, there is $\eta>0$ such that whenever $\norm{z}<\eta$,
\[
  \norm{x+z}=\varphi(x+z)
  \quad\text{and}\quad
  \norm{y+z}=\varphi(y+z).
\]
Consequently, for all $\norm{z}<\eta$,
\[
  \norm{x+z}-\norm{y+z}
  =
  \varphi(x+z)-\varphi(y+z)
  =
  \varphi(x)-\varphi(y)
  =
  0.
\]
The distinct unit vectors $x,y$ therefore violate property $(S)$ with
$a=\eta$.

If $|M|=2$, then $\Free(M)$ is one-dimensional. Its unit sphere consists
of two opposite points, say $\{\xi,-\xi\}$. For every $0<d<2$ and $a>0$,
the only admissible pair is $\xi,-\xi$, and choosing
$z=t\xi$ with $0<t<\min\{a,1\}$ gives
\[
  |\norm{\xi+z}-\norm{-\xi+z}|=2t>0.
\]
Hence the one-dimensional space has property $(S)$; moreover the
quantitative modulus is positive for each $d,a$, so it has uniform
property $(S)$.
\end{proof}

\section*{Verification Notes}

No computational verifier was used for this packet. The proof is a
finite-dimensional convex-geometric argument. The main review points are:

\begin{itemize}
\item the standard molecule description of the unit ball of $\Free(M)$
for finite $M$;
\item the fact that a polyhedral norm is locally represented by a single
facet functional at relative interior points of a facet;
\item the one-dimensional endpoint case $|M|=2$.
\end{itemize}

\section*{Limitations and Novelty Check}

This proves only the finite pointed metric-space subcase of the
Lipschitz-free open problem. It does not characterize infinite metric
spaces. The result is likely folklore-level once the finite free-space
unit ball is recalled, but it is a clean deterministic subcase suitable
for the pipeline's first Banach-space smoke test.

\end{document}
